\section{Phase mixing as decorrelation of slices}\label{sec:decorrelation}

By theorem~\ref{thm:slice}, one compressive branch is a continuum of passive scalars $a_u$, each carried by its own divergence-free velocity $\boldsymbol U_u=\uE+su\,\bhat$ \eqref{eq:Uu}. Without collisions, sources or dissipation of $G$, the energy of every slice is conserved (section~\ref{sec:invariants}). A Hermite moment $g_m$ is a fixed linear combination of slices, so free energy can move between Hermite moments only through changes in the correlations between different slices. This section makes the statement precise. We derive the evolution law of the two-slice correlation (section~\ref{dec:sec-twoslice}), write the Hermite flux and the Hermite spectrum in terms of it (sections~\ref{dec:sec-flux} and~\ref{dec:sec-dictionary}), examine the split of the flux into phase-mixing and phase-unmixing parts that uses the Eulerian wavenumber $k_z$ (section~\ref{dec:sec-tilt}), and show that the invariants of section~\ref{sec:invariants} do not constrain the transfer (section~\ref{dec:sec-fjortoft}).

Throughout this section $\langle\cdot\rangle$ is the average over the periodic box,
\begin{equation*}
  C(u,u')=\langle a_u\,a_{u'}\rangle
\end{equation*}
is the two-slice correlation, and $W(m)=p_m\langle g_m^2\rangle$ is the Hermite spectrum of the free energy. The diagonal $C(u,u)$ is the slice energy $\mathcal E(u)$ of \eqref{eq:Qphi}. Since $\fe{G}{G}=\sum_mp_mg_m^2$ and, by theorem~\ref{thm:spectrum}, $\fe{G}{G}=\int a_u^2\,\dd u$, the total $\sum_mW(m)=\langle\fe{G}{G}\rangle=\int C(u,u)\,\dd u$ is $2W^\pm/n_iT_i$ in the normalisation of \eqref{eq:WGH}. For the Fourier coefficients of a mode, products such as $a_ua_{u'}$ are replaced by $a_ua_{u'}^*$. In the truncated hierarchy of section~\ref{sec:truncated} the slices are the unit-weight amplitudes $a_j$ at the eigenvalues $u_j$ of the ladder matrix, $j=0,\dots,M$, and the two-slice correlation is the matrix $C_{jl}=\langle a_ja_l\rangle$.

\subsection{The two-slice law}\label{dec:sec-twoslice}

\begin{proposition}[Two-slice law]\label{prop:twoslice}
Assume the hypotheses of theorem~\ref{thm:slice}, with no collision operator, source or dissipation acting on $G$, and assume that the slice amplitudes $a_u(\boldsymbol r,t)$ are continuously differentiable in $\boldsymbol r$ and $t$, with derivatives that are continuous in $u$. Then, for all $u$ and $u'$,
\begin{equation}\label{eq:twoslice}
  \frac{\dd C(u,u')}{\dd t}=-s\,(u-u')\,\big\langle a_{u'}\,\gradpar a_u\big\rangle .
\end{equation}
In particular, $C(u,u)$ is constant for every $u$, and the $\boldsymbol E\times\boldsymbol B$ flow does not enter.
\end{proposition}

\begin{proof}
Since $\partial_z+\{\Psi,\cdot\}=\gradpar$, the slice equation \eqref{eq:slice} reads $\partial_ta_u=-\{\Phi,a_u\}-su\,\gradpar a_u$. The bracket and $\gradpar$ obey the Leibniz rule, so
\begin{equation*}
  \partial_t(a_ua_{u'})=-\{\Phi,a_ua_{u'}\}-su'\,\gradpar(a_ua_{u'})-s\,(u-u')\,a_{u'}\gradpar a_u .
\end{equation*}
For every continuously differentiable $f$, $\{\Phi,f\}=\nabla\cdot(f\,\uE)$ and $\gradpar f=\nabla\cdot(f\,\bhat)$, because $\nabla\cdot\uE=\nabla\cdot\bhat=0$. The first two terms are therefore divergences, and their averages over the periodic box vanish.
\end{proof}

The $\boldsymbol E\times\boldsymbol B$ flow carries the two slices of a pair together and cannot change their correlation; only the difference of their velocities, $\boldsymbol U_u-\boldsymbol U_{u'}=s(u-u')\,\bhat$, enters \eqref{eq:twoslice}. The diagonal $C(u,u)$ is the energy of slice $u$, whose conservation gives the invariants of section~\ref{sec:invariants}. The correlation between two different slices changes only through their relative streaming along the perturbed field, at a rate set by the difference of their speeds. This gives phase mixing and phase unmixing a definition that does not refer to any wavenumber: they are the decorrelation and the recorrelation of slices along $\gradpar$. Under the hypotheses \eqref{eq:hyp} of theorem~\ref{thm:frozen}, which include ideal induction and exclude spatial dissipation of $G$, the law can be integrated. Section~\ref{sec:lagrangian} shows that $C(u,u';t)$ is then the correlation of the initial slices at the field-aligned separation $s(u'-u)t$ along the field lines frozen at $t=0$ \eqref{eq:frozencorr}, so that phase mixing is the decay of the initial correlations with that separation.

The proposition needs no assumption on the Alfv\'enic fields beyond smoothness; they may be forced, dissipated or prescribed. Operators that act on $G$ itself add terms to \eqref{eq:twoslice}. Collisions and sources enter through their slice form \eqref{eq:sliceforcing}, and a velocity-independent dissipation, such as perpendicular hyperdiffusion applied to every $g_m$, acts on each slice separately and changes $C(u,u)$ as well. The same algebra holds for the truncated hierarchy (theorem~\ref{thm:trunc}), with $C_{jl}$ and the nodes $u_j$, $u_l$ in place of $C(u,u')$ and $u$, $u'$. In a pseudo-spectral computation the step from the pointwise identity to \eqref{eq:twoslice} uses the exact antisymmetry of the discrete bracket under the volume average, so it requires $\Phi$, $\Psi$ and $g$ to lie inside the dealiasing band; the slice equation itself does not. We checked the truncated form of \eqref{eq:twoslice} to round-off against GANDALF's right-hand side \citep{Kanekar2025}, for $M=8$, 9 and 12 and for $M=64$ and 65, $\Lam=\pm\sqrt5$ and $3$, and both closures, in states in which the $\boldsymbol E\times\boldsymbol B$ term of a single Fourier mode is as large as its total rate of change. The two sides agree to a relative accuracy of $5\times10^{-15}$, and the $\boldsymbol E\times\boldsymbol B$ contribution to the averaged law vanishes to $2\times10^{-16}$.

The average over the box is essential. With $f(\boldsymbol r)=\sum_{\boldsymbol k}f(\boldsymbol k)\,e^{\ii\boldsymbol k\cdot\boldsymbol r}$ and $\{f,h\}_{\boldsymbol k}$ the Fourier coefficient of the bracket, the correlation of a single Fourier mode, $C_{\boldsymbol k}(u,u')=a_u(\boldsymbol k)\,a_{u'}(\boldsymbol k)^*$, obeys
\begin{equation}\label{dec:pairk}
  \frac{\dd C_{\boldsymbol k}(u,u')}{\dd t}=-\ii sk_z\,(u-u')\,C_{\boldsymbol k}(u,u')+N^\Phi_{\boldsymbol k}(u,u')+N^\Psi_{\boldsymbol k}(u,u'),
\end{equation}
where
\begin{equation*}
  N^\Phi_{\boldsymbol k}=-\{\Phi,a_u\}_{\boldsymbol k}\,a_{u'}(\boldsymbol k)^*-a_u(\boldsymbol k)\,\{\Phi,a_{u'}\}_{\boldsymbol k}^*,\qquad
  N^\Psi_{\boldsymbol k}=-su\,\{\Psi,a_u\}_{\boldsymbol k}\,a_{u'}(\boldsymbol k)^*-su'\,a_u(\boldsymbol k)\,\{\Psi,a_{u'}\}_{\boldsymbol k}^* .
\end{equation*}
The first term rotates the phase of $C_{\boldsymbol k}$ at the Eulerian rate $sk_z(u-u')$. The $\boldsymbol E\times\boldsymbol B$ term $N^\Phi_{\boldsymbol k}$ moves correlation between wavenumbers and sums to zero over $\boldsymbol k$. The term $N^\Psi_{\boldsymbol k}$ sums to $-s(u-u')\langle a_{u'}\{\Psi,a_u\}\rangle$ and completes $\partial_z$ to $\gradpar$ in \eqref{eq:twoslice}. Mode by mode, therefore, the $\boldsymbol E\times\boldsymbol B$ advection does not drop out, and the linear dephasing is measured with $k_z$ rather than along the field. In the states of our checks the largest $\boldsymbol E\times\boldsymbol B$ term of a single mode was 1.0--1.2 times the largest total rate of change of a mode, while the sum of these terms over modes vanished to round-off.

In electrostatic drift kinetics with a straight guide field \citep{Schekochihin2016,Parker2016}, and in the one-dimensional linear model of \citet{Kanekar2015}, phase mixing and phase unmixing are separated mode by mode with the sign of the parallel wavenumber. The Hermite moments of each Fourier mode are split into a phase-mixing part, which propagates to larger $m$, and a phase-unmixing part, which propagates to smaller $m$, and the Hermite flux of the mode is proportional to the difference of their energies. In KRMHD the field lines move. \citet{Meyrand2019} noted that the sign would then have to be taken with respect to the perturbed field, that they did not know how to make such a decomposition rigorous, and that they therefore did not use one. The two-slice law describes phase mixing along the perturbed field without reference to a wavenumber. Section~\ref{dec:sec-tilt} examines what is reported when the sign of $k_z$ is used instead.

\subsection{The Hermite flux as a Christoffel--Darboux transform}\label{dec:sec-flux}

We first write the Hermite moments in slice variables. Let $r_n(u)$ be the $n$th Hermite coefficient of the Case--Van Kampen mode $F_u$ (section~\ref{sec:slices}). It is a polynomial of degree $n$, with $r_0=1$ and $r_1=\sqrt2\,u$, and the slice amplitude of the Hermite basis function $e_n=H_nF_M/\sqrt{2^nn!}$ is $\sqrt\varrho\,p_nr_n$. Since $\fe{e_n}{G}=p_ng_n$, the unitarity of $G\mapsto a$ (theorem~\ref{thm:spectrum}) gives
\begin{equation}\label{dec:gn}
  g_n=\int\sqrt{\varrho(u)}\;r_n(u)\,a_u\,\dd u=\frac{1}{\sqrt{p_n}}\int\phi_n(u)\,a_u\,\dd u,\qquad \phi_n=\sqrt{p_n\varrho}\;r_n ,
\end{equation}
where the $\phi_n$ are the orthonormal basis of $L^2(\mathbb R,\dd u)$ of section~\ref{sec:slices}. Let $E_{\le m}=\sum_{n\le m}W(n)$ be the free energy in the Hermite moments $n\le m$. By \eqref{dec:gn},
\begin{equation}\label{dec:kernel}
\begin{aligned}
  E_{\le m}&=\iint K_m(u,u')\,C(u,u')\,\dd u\,\dd u',\\
  K_m(u,u')&=\sum_{n\le m}\phi_n(u)\,\phi_n(u')=\sqrt{\varrho(u)\varrho(u')}\,\sum_{n\le m}p_n\,r_n(u)\,r_n(u') ,
\end{aligned}
\end{equation}
so $K_m$ is the kernel of the orthogonal projection, in $L^2(\mathbb R,\dd u)$, onto the span of $\phi_0,\dots,\phi_m$, that is, onto the Hermite moments $n\le m$.

Let $b_m=p_mA_{m,m+1}$, so that $b_0=(1-1/\Lam)/\sqrt2$ and $b_m=\sqrt{(m+1)/2}$ for $m\ge1$; the symmetry of $\Pw A$ is the statement $b_m=p_{m+1}A_{m+1,m}$. Under the hypotheses of proposition~\ref{prop:twoslice}, the Hermite coefficients obey $\partial_tg+\{\Phi,g\}+s\gradpar(Ag)=0$, the Hermite form of \eqref{eq:branch}. The $\boldsymbol E\times\boldsymbol B$ term does not change $E_{\le m}$, because the bracket is antisymmetric under the volume average. The streaming term gives $\dd E_{\le m}/\dd t=-2s\sum_{n\le m}\sum_{n'}(\Pw A)_{nn'}\langle g_n\gradpar g_{n'}\rangle$, in which $\Pw A$ is symmetric and $\langle g_n\gradpar g_{n'}\rangle$ is antisymmetric in $(n,n')$. The terms with $n,n'\le m$ cancel in pairs, and only the coupling between $m$ and $m+1$ remains. Differentiating \eqref{dec:kernel} instead, with the two-slice law, gives the same rate in slice variables. Hence
\begin{equation}\label{eq:fluxCD}
\begin{aligned}
  \frac{\dd E_{\le m}}{\dd t}=-\Pi_{m+1/2},\qquad
  \Pi_{m+1/2}&=s\,b_m\big\langle g_m\gradpar g_{m+1}-g_{m+1}\gradpar g_m\big\rangle\\
  &=-\iint K_m(u,u')\,\frac{\partial C(u,u')}{\partial t}\,\dd u\,\dd u' .
\end{aligned}
\end{equation}
The flux $\Pi_{m+1/2}$ is the free energy per unit time and per unit volume that passes from moment $m$ to moment $m+1$. The first line is the Hermite flux of \citet{Schekochihin2016} and \citet{AdkinsSchekochihin2018} written with the full field-aligned gradient $\gradpar=\partial_z+\{\Psi,\cdot\}$; \citet[eq.~9]{Meyrand2019} use the field-aligned gradient in the same way, in a flux that they also resolve in $k_\perp$.

The two lines of \eqref{eq:fluxCD} are related term by term by a Christoffel--Darboux identity. The operator $L$ acts on the Hermite basis as $Le_n=A_{n-1,n}e_{n-1}+A_{n+1,n}e_{n+1}$, and $\Tr[LG]=u\,\Tr[G]$. Taking the transform of $Le_n$ therefore gives
\begin{equation*}
  u\,p_n\,r_n=p_{n-1}A_{n-1,n}\,r_{n-1}+p_{n+1}A_{n+1,n}\,r_{n+1},
\end{equation*}
and the symmetry of $\Pw A$ turns this into the three-term recurrence
\begin{equation*}
  u\,r_n(u)=A_{n,n+1}\,r_{n+1}(u)+A_{n,n-1}\,r_{n-1}(u),\qquad r_{-1}=0 .
\end{equation*}
Multiply the recurrence at $u$ by $p_nr_n(u')$, subtract the same expression with $u$ and $u'$ exchanged, and sum over $n\le m$. Because $p_nA_{n,n+1}=b_n$ and $p_nA_{n,n-1}=b_{n-1}$, the sum telescopes, and
\begin{equation}\label{dec:CD}
  (u-u')\,K_m(u,u')=b_m\sqrt{\varrho(u)\varrho(u')}\,\big[r_{m+1}(u)\,r_m(u')-r_m(u)\,r_{m+1}(u')\big] .
\end{equation}
With the two-slice law, the second line of \eqref{eq:fluxCD} becomes
\begin{equation*}
  \Pi_{m+1/2}=s\iint(u-u')\,K_m(u,u')\,\big\langle a_{u'}\gradpar a_u\big\rangle\,\dd u\,\dd u' ,
\end{equation*}
and \eqref{dec:CD} together with \eqref{dec:gn} turns this integral into the first line. The pairs $u=u'$ carry no flux, because $\partial_tC(u,u)=0$. The free energy that passes through every $m$ is carried entirely by the correlations between different slices.

The truncated hierarchy has the same structure, with sums over the nodes. Let $r_j$ be the right eigenvectors of the truncated ladder matrix, normalised by $(r_j)_0=1$, let $n_j=r_j^{\mathsf T}\Pw r_j$, and let $R_{mj}=(r_j)_m/\sqrt{n_j}=r_m(u_j)/\sqrt{n_j}$, so that $R^{\mathsf T}\Pw R=I$, $g_m=\sum_jR_{mj}a_j$ and $a_j=\sum_mp_mR_{mj}g_m$. Then
\begin{equation}\label{dec:fluxtrunc}
\begin{aligned}
  &\Pi_{m+1/2}=-\sum_{j,l}K_m(j,l)\,\dot C_{jl},\qquad
  K_m(j,l)=\sum_{n\le m}p_n\,R_{nj}R_{nl},\\
  &b_m\big(R_{mj}R_{m+1,l}-R_{m+1,j}R_{ml}\big)=(u_l-u_j)\,K_m(j,l),
\end{aligned}
\end{equation}
with $\dot C_{jl}=-s(u_j-u_l)\langle a_l\gradpar a_j\rangle$. The last relation is \eqref{dec:CD} at the nodes, and it holds for every $m<M$. The copy closure changes only the last row of the ladder matrix and the weight of $g_M$ (appendix~\ref{app:truncated}); with the eigenvectors normalised by its weight, \eqref{dec:fluxtrunc} holds unchanged for $m<M$. We checked the budget \eqref{eq:fluxCD}, the identity and the slice form of the flux to round-off with GANDALF's ladder matrix and right-hand side, for $M=8$, 9 and 12 and for $M=64$ and 65, $\Lam=\pm\sqrt5$ and $3$, and both closures, in states with strong brackets; the relative errors are below $5\times10^{-14}$.

At $\Lam=\infty$, $\varrho=F_M$, $p_n=1$ and $\phi_n=\psi_n$, the orthonormal Hermite functions. The kernel $K_m$ is then the classical Christoffel--Darboux kernel of the Hermite functions, and for $|u|,|u'|\ll\sqrt{2m}$ it approaches the sine kernel,
\begin{equation*}
  K_m(u,u')\simeq\frac{\sin[\kappa_m(u-u')]}{\pi\,(u-u')},\qquad \kappa_m=\sqrt{2m+1};
\end{equation*}
on $|u|,|u'|<1.5$ the largest deviation is 0.8\%, 0.4\% and 0.2\% of the maximum of the kernel at $m=64$, 128 and 256. The factor $u-u'$ in the pair form of the flux cancels the denominator, so that the pair weight $(u-u')K_m$ approaches $\sin[\kappa_m(u-u')]/\pi$; at separations of order one this approximation is less accurate than that of the kernel. For slices concentrated at $|u|\ll\sqrt{2m}$,
\begin{equation}\label{eq:fluxsine}
  \Pi_{m+1/2}\simeq\frac{s}{\pi}\iint\sin[\kappa_m(u-u')]\,\big\langle a_{u'}\gradpar a_u\big\rangle\,\dd u\,\dd u'
  =\frac{s}{\pi}\int\sin(\kappa_m\delta)\,\Theta(\delta)\,\dd\delta ,
\end{equation}
where $\Theta(\delta)=\int\langle a_{\bar u-\delta/2}\,\gradpar a_{\bar u+\delta/2}\rangle\,\dd\bar u$ is the pair correlation at the velocity separation $\delta$. It is odd in $\delta$, because $\gradpar$ is antisymmetric under the volume average. In the truncated system at $\Lam=\infty$ a pair of slices $(j,l)$ enters with the weight $s\sqrt{\Delta_j\Delta_l}\,\sin[\kappa_m(u_j-u_l)]/\pi$ multiplying $\langle a_l\gradpar a_j\rangle$, where $\Delta_j$ is the local spacing of the nodes.

At $\Lam=\infty$, up to the factor $s/\pi$, the flux through $m$ is therefore the sine transform of the antisymmetric pair correlation $\Theta$ at the velocity-space wavenumber $\kappa_m$. The Hermite index selects a Fourier component of the dephasing pattern, not a band of slice separations. Like $\sin(\kappa_m\delta)/\pi$, the exact pair weight does not decay with $|\delta|$: at $m=128$ its envelope stays between 0.29 and 0.37, close to $1/\pi\approx0.32$, out to $|\delta|=6$. In the random states of our checks, at $\Lam=-\sqrt5$ and 3, the pairs closer than $2\pi/\kappa_m$ carried only 6--22\% of $\sum_{j,l}|K_m(j,l)\,\dot C_{jl}|$ at $m=8$, 16 and 32. In free streaming at $\Lam=\infty$ from a density perturbation at one wavenumber $k_z$, the pair correlation $\langle a_{u'}\gradpar a_u\rangle$ is proportional to $[F_M(u)F_M(u')]^{1/2}\sin[sk_zt\,(u-u')]$, a dephasing pattern whose velocity-space wavenumber $s|k_z|t$ grows linearly in time, and the flux through $m$ is largest at $s|k_z|t=\kappa_m$ (appendix~\ref{app:tilt}).

Equation~\eqref{eq:fluxsine} is a statement about $\Lam=\infty$, that is, about the Laguerre sectors with $l\ge2$ and about free streaming. At finite $\Lam$ the kernel deviates from the sine kernel, on the same range of $u$, by 11.8\%, 8.5\% and 6.1\% of its maximum at $m=64$, 128 and 256 for $\Lam=\sqrt5$, and by 4.6\%, 3.3\% and 2.4\% for $\Lam=-\sqrt5$. These relative deviations decrease only because the maximum of the kernel grows as $\kappa_m/\pi$. The pair weight $(u-u')K_m$ does not approach $\sin[\kappa_m(u-u')]/\pi$ at finite $\Lam$, and its deviation from it remains of order $1/\pi$ as $m$ grows.

The field-line part of the flux, which comes from $\{\Psi,\cdot\}$ in $\gradpar$, is not small. Keeping only $\partial_z$ in \eqref{eq:fluxCD} gives the $k_z$-based flux, a sum over Fourier modes of
\begin{equation}\label{dec:fluxz}
  \Pi^z_{m+1/2}(\boldsymbol k)=-2s\,b_m\,k_z\,\mathrm{Im}\big[g_m(\boldsymbol k)^*g_{m+1}(\boldsymbol k)\big].
\end{equation}
GANDALF's diagnostic \texttt{hermite\_flux} returns $\Pi^z_{m+1/2}(\boldsymbol k)/(s\,p_m)$: it keeps only $\partial_z$, and it omits the factor $s$ and, at $m=0$, the factor $p_0$. In the strongly nonlinear random states of our checks the field-line part, $\Pi_{m+1/2}-\sum_{\boldsymbol k}\Pi^z_{m+1/2}(\boldsymbol k)$, was between 29\% and 253\% of $\Pi_{m+1/2}$ in magnitude, depending on $m$. In slice variables the $k_z$-based flux of a mode is the Christoffel--Darboux filter of \eqref{dec:fluxtrunc} applied to the Eulerian linear term of \eqref{dec:pairk} alone,
\begin{equation*}
  \Pi^z_{m+1/2}(\boldsymbol k)=-\sum_{j,l}K_m(j,l)\,\mathrm{Re}\big[-\ii sk_z\,(u_j-u_l)\,C_{\boldsymbol k}(j,l)\big],
\end{equation*}
with $C_{\boldsymbol k}(j,l)=a_j(\boldsymbol k)\,a_l(\boldsymbol k)^*$.

\subsection{The Hermite spectrum and the two-slice correlation}\label{dec:sec-dictionary}

By \eqref{dec:gn}, the Hermite spectrum is a quadratic form in the two-slice correlation:
\begin{equation}\label{eq:dictionary}
  W(m)=\iint C(u,u')\,\phi_m(u)\,\phi_m(u')\,\dd u\,\dd u' .
\end{equation}
At $\Lam=\infty$ the slices are the values of the distribution function, $a_u=G(u)/\sqrt{F_M(u)}$, and $\phi_m=\psi_m$. We call the slices $\delta$-correlated if $C(u,u')=e(u)\,\delta(u-u')$ with a density $e\ge0$. They are then completely decorrelated in velocity, and \eqref{eq:dictionary} gives $W(m)=\int e\,\phi_m^2\,\dd u$. This spectrum decays only as $m^{-1/2}$ (see \eqref{eq:prefactor} below), so $\sum_mW(m)$ diverges: $\delta$-correlated slices are an idealisation of states with velocity structure on arbitrarily fine scales, such as the collisionless forced state discussed below, whose total free energy diverges in the same way \citep{Kanekar2015}. For $e=F_M$ at $\Lam=\infty$,
\begin{equation}\label{eq:prefactor}
  \int F_M\,\psi_m^2\,\dd u=\frac{\Gamma(m+\frac12)}{\sqrt2\,\pi\,m!}=\frac{1}{\pi\sqrt{2m}}\Big(1-\frac{1}{8m}+O(m^{-2})\Big),
\end{equation}
where $\Gamma$ is Euler's gamma function.

\citet{Kanekar2015} studied a one-dimensional linear kinetic model that has the form of one branch, with $\Lam=-1/\alpha$ in terms of their parameter $\alpha$, and whose density moment is driven by white noise (their $\chi_k$, with $\langle\chi_k(t)\chi_k^*(t')\rangle=\epsilon_k\,\delta(t-t')$). Their eq.~(4.10), derived without the collision term of their model, gives the exact steady-state Hermite spectrum of a Fourier mode for $m\ge1$, which in our variables reads
\begin{equation*}
  W(m)=\frac{(1-1/\Lam)^2\,\epsilon_k}{2\pi\,s|k_z|}\,\frac{1}{2^mm!}\int\frac{|Z^{(m)}(\zeta)|^2}{|\eps(\zeta+\ii0)|^2}\,\dd\zeta ,
\end{equation*}
where $Z^{(m)}$ is the $m$th derivative of the plasma dispersion function, $\eps$ is the dielectric function \eqref{eq:eps}, and their $|k|$ is our $s|k_z|$. At $\Lam=\infty$ the dielectric factor drops out. On the real axis $\mathrm{Im}\,Z^{(m)}=(-1)^m\pi H_mF_M$ and $\mathrm{Re}\,Z^{(m)}=(-1)^m\pi\,\mathsf H[H_mF_M]$, with $\mathsf H$ the Hilbert transform of section~\ref{sec:slices}. Because $\mathsf H$ preserves the $L^2$ norm,
\begin{equation*}
  \frac{1}{2^mm!}\int|Z^{(m)}(\zeta)|^2\,\dd\zeta=\frac{2\pi^2}{2^mm!}\int H_m^2F_M^2\,\dd\zeta=\frac{\sqrt2\,\pi\,\Gamma(m+\frac12)}{m!},
\end{equation*}
and their spectrum becomes
\begin{equation}\label{dec:kanekar}
  W(m)=\frac{\pi\epsilon_k}{s|k_z|}\,\frac{\Gamma(m+\frac12)}{\sqrt2\,\pi\,m!} .
\end{equation}
This is \eqref{eq:prefactor} with the density $e=(\pi\epsilon_k/s|k_z|)F_M$: the forced linear spectrum is the spectrum of $\delta$-correlated slices. The closed form is not written in their paper, whose eq.~(4.14) is its large-$m$ limit, and the slice argument covers $m=0$ as well.

The reason for the $\delta$-correlation is dephasing. Density forcing loads every slice with the same spatial pattern \eqref{eq:sliceforcing}. Each slice then streams at its own speed $su$, so two slices loaded together dephase at the rate $s|k_z(u-u')|$. In the statistically steady state, which superposes all past injections, the real part of the mean of $a_u(\boldsymbol k)\,a_{u'}(\boldsymbol k)^*$ therefore vanishes for $u\ne u'$, so that the real-space correlation $C(u,u')$ is $\delta$-correlated. The imaginary part, which is antisymmetric in $(u,u')$ and does not enter \eqref{eq:dictionary}, does not vanish; through $\langle a_{u'}\gradpar a_u\rangle$ it carries the flux $\epsilon_W=(1-1/\Lam)\,\epsilon_k$, the injection rate of $\sum_mW(m)$, through every $m$.

At finite $\Lam$ the argument is the same, and the real-space correlation is again $\delta$-correlated, now with $e=\pi(1-1/\Lam)^2\epsilon_k\,\varrho/(s|k_z|)$, which reproduces their eq.~(4.10) exactly; but $e$ and $\phi_m$ then differ from $F_M$ and $\psi_m$, and the closed form \eqref{dec:kanekar} holds only asymptotically. At $\Lam=\sqrt5$ the exact spectrum is 4.7\% below $(1-1/\Lam)$ times \eqref{dec:kanekar} at $m=1$. At large $m$, $W(m)\to\epsilon_W/(s|k_z|\sqrt{2m})$. This level does not depend on $\Lam$, because of their eq.~(4.15), which in our variables reads $\int F_M\,|\eps(\zeta+\ii0)|^{-2}\,\dd\zeta=1/(1-1/\Lam)$. In slice variables (4.15) is the normalisation $\int\varrho\,\dd u=1/p_0$ of the spectral density (section~\ref{sec:slices}), by which the slice energy injected by density forcing, at the rate $(1-1/\Lam)^2\epsilon_k\,\varrho(u)$ per unit $u$, adds up to $\epsilon_W$.

Correlated slices give other spectra. Let $C(u,u')=[e(u)e(u')]^{1/2}\,c(u-u')$ at $\Lam=\infty$, with $c$ even, $e$ varying slowly on the correlation length of $c$, and $\hat c(\kappa)=\int c(\delta)\cos(\kappa\delta)\,\dd\delta$ the spectrum of the slice correlation in the velocity-space wavenumber $\kappa$. For $|u|,|u'|\ll\sqrt{2m}$, $\psi_m(u)\psi_m(u')\approx\{\cos[\sqrt{2m}(u-u')]+\cos[\sqrt{2m}(u+u')-m\pi]\}/(\pi\sqrt{2m})$, and the second term averages out against the slowly varying $e$. Then \eqref{eq:dictionary} gives
\begin{equation}\label{dec:powerlaw}
  W(m)\simeq\Big[\int e\,\dd u\Big]\,\frac{\hat c(\sqrt{2m})}{\pi\sqrt{2m}} ,
\end{equation}
so the Hermite index $m$ probes the slice correlation at the velocity-space wavenumber $\sqrt{2m}$. A power-law spectrum $\hat c\propto\kappa^{-p}$ gives $W(m)\propto m^{-(p+1)/2}$, and the $\delta$-correlated case is $p=0$. For $1<p<3$ the same correlation has the second-order structure function $\langle(a_{u+\Delta}-a_u)^2\rangle\propto\Delta^{p-1}$ at small $\Delta$. Exponentially correlated slices, $c=e^{-|\delta|/\ell}$, have $p=2$ for $m\gg1/(2\ell^2)$; for $\ell=0.3$ and $e=F_M$, direct quadrature of \eqref{eq:dictionary} agrees with \eqref{dec:powerlaw} to 0.2\% at $m=200$, and the local slope of $W(m)$ between $m=100$ and 200 is $-1.463$ against $-1.462$ predicted. The correspondence is restricted to power-law spectra. If $\hat c$ decays faster than any power, the terms neglected in \eqref{dec:powerlaw} dominate at large $m$: for a Gaussian correlation $c=e^{-\delta^2/2\ell^2}$ with $\ell=0.3$, the exact $W(m)$ exceeds \eqref{dec:powerlaw} by factors of 1.16, 1.40 and 2.06 at $m=50$, 100 and 200. Figure~\ref{fig:dictionary} illustrates the dictionary.

That at large $m$ the Hermite index acts as a velocity-space wavenumber $\sqrt{2m}$, so that the Hermite transform is effectively a Fourier transform in velocity, is the duality noted by \citet{Schekochihin2016} and \citet{AdkinsSchekochihin2018}; given it, the relation between the exponents is elementary. \citet[note~81]{Eyink2018} related the $m^{-1/2}$ spectrum to a velocity-space regularity exponent of $-1/2$, and \citet{Nastac2024} describe linear phase mixing as producing a spectrum that is flat in the velocity-space wavenumber up to a collisional cutoff, which is the case $p=0$. What the slice picture adds is the representation \eqref{eq:dictionary} of the Hermite spectrum through the two-slice correlation, the closed form \eqref{dec:kanekar} of the forced linear spectrum, and the reading of the $m^{-1/2}$ spectrum as complete velocity-space decoherence of the slices in their equal-time correlation.

\subsection{The \texorpdfstring{$k_z$}{kz}-based decomposition under field-line tilt}\label{dec:sec-tilt}

Some diagnostics of phase mixing in KRMHD simulations use the Eulerian wavenumber $k_z$ in place of the field-aligned one. The $k_z$-based flux \eqref{dec:fluxz} keeps only $\partial_z$ in \eqref{eq:fluxCD}, and the split into phase-mixing and phase-unmixing parts is made mode by mode with $\mathrm{sgn}\,k_z$. The KRMHD diagnostics of \citet[chapter~6]{KanekarThesis2015} work in this way, and so do GANDALF's diagnostics \texttt{phase\_mixing\_energy} and \texttt{phase\_unmixing\_energy}, which assign $|g_m(\boldsymbol k)|^2$ of each Fourier mode to phase mixing or to unmixing according to the sign of its $k_z$-based flux. \citet{Meyrand2019} did not use such a split (section~\ref{dec:sec-twoslice}). To see what these diagnostics report when the field lines are tilted, we consider streaming along a uniformly tilted static field, $\gradpar=\partial_z+b_0\partial_x$, with no flow. Each Fourier mode then evolves as in a uniform field with the field-aligned wavenumber $k_\parallel=k_z+b_0k_x$. A uniform tilt cannot be produced by a periodic $\Psi$, but it is the local form of the tilt that a large-scale $\boldsymbol b_\perp$ imposes on structure of much smaller scale, and the frozen-in shear of section~\ref{sec:lagrangian} produces one exactly, with a tilt that grows in time. Appendix~\ref{app:tilt} gives the derivations.

On each mode $\Pi^z_{m+1/2}(\boldsymbol k)=(k_z/k_\parallel)\,\Pi_{m+1/2}(\boldsymbol k)$. On the modes with $k_zk_\parallel<0$ the $k_z$-based flux therefore has the wrong sign, and the $\mathrm{sgn}\,k_z$ split exchanges the phase-mixing and phase-unmixing parts. Free streaming from a density perturbation is pure phase mixing, with a field-aligned flux that is non-negative on every mode and at every $m$ and $t$, yet a mode with $k_z<0<k_\parallel$ is reported as pure unmixing. Consider an ensemble of modes with a Gaussian spectrum in $k_z$ and $k_x$, with initial conditions that share one velocity profile, and define the effective tilt angle $\alpha$ by letting $\tan\alpha$ be the ratio of the r.m.s.\ value of $b_0k_x$ to that of $k_z$. The reported net flux and the spectral weight of the misclassified modes are
\begin{equation}\label{eq:tiltcos}
  \Pi^z_{m+1/2}(t)=\cos^2\alpha\;\Pi_{m+1/2}(t),\qquad
  \frac{\text{weight of the modes with }k_zk_\parallel<0}{\text{total weight}}=\frac{\alpha}{\pi},
\end{equation}
for every $m$ and $t$, every $\Lam$, and with any $m$-diagonal collision operator. For pure phase mixing the reported backward flux, the sum over modes of the magnitudes of the negative $k_z$-based fluxes, is positive although the true backward flux is zero. At fixed $m$ its share of the gross reported flux tends to $1/2$ at late times (table~\ref{appE:tab}); the gross reported flux then exceeds the true flux by a growing factor, while the net reported flux stays at $\cos^2\alpha$ times the true flux. On an exact solution of this kind with $\alpha=45^\circ$, GANDALF's \texttt{phase\_unmixing\_energy} assigns 33\% of the free energy at $m=1$ to unmixing by $s\sigma_\parallel t=6$, where $\sigma_\parallel$ is the r.m.s.\ field-aligned wavenumber of the ensemble (appendix~\ref{app:tilt}; see also figure~\ref{fig:tilt}). In critically balanced turbulence $b_\perp\sim k_\parallel/k_\perp$ at each scale, so $b_0k_x$ is of the same order as $k_z$ and $\tan\alpha$ is of order unity. The model then predicts that the $k_z$-based split assigns an order-unity fraction of phase mixing to unmixing; we have not measured this in a turbulence simulation.

The statement that the true backward flux vanishes holds for free streaming, and we found it to hold for density initial data at $\Lam=\sqrt5$ (the branch $\Lam^+$ at $\beta_i=\tau=Z=1$) for $m<2600$ and $s|k_\parallel|t\le60$. It is not general: with density data at $\Lam=-\sqrt5$, and with data in $g_1$ at $\Lam=\sqrt5$, the field-aligned flux through $m=1/2$ becomes negative for a time, by up to 2.2\% and 36\% of the largest flux at that time (appendix~\ref{app:tilt}). The relation \eqref{eq:tiltcos} for the net flux does not depend on this premise. The misattribution is a property of the diagnostics and not of the dynamics. The field-aligned flux \eqref{eq:fluxCD} and the two-slice law involve no wavenumber, and, under the hypotheses \eqref{eq:hyp} and without collisions, no weight moves between field-aligned wavenumbers, although the Eulerian $k_z$ of a Fourier mode can change sign (corollary~\ref{cor:kpar}).

\subsection{No Fj{\o}rtoft-type constraint}\label{dec:sec-fjortoft}

\citet[section~4.1]{Chandran2026} discuss cascade directions with the generalisation by \citet{AlexakisBiferale2018} of the argument of \citet{Fjortoft1953} to invariants that are not sign-definite: if two invariants $\mathcal I_1$ and $\mathcal I_2$ have spectra with $E_{\mathcal I_2}\ge c\,k^n|E_{\mathcal I_1}|$ for positive constants $c$ and $n$ (their condition 4.1), then $\mathcal I_1$ cannot cascade to large $k$ at a constant rate. They apply the same argument with $k$ replaced by the Hermite index $m$. In their appendix~B they note that $\Gamma$, whose Hermite form carries the factor $\sqrt{m+1}$, is not sign-definite, so that $W$ and $\Gamma$ do not satisfy the condition, and they expect both to cascade to large $m$. In their conclusion they ask more generally whether the conservation of helicity leads to a bottleneck in velocity space. The slice structure completes their argument for $\Gamma$ and extends it to every invariant of collisionless KRMHD that depends on $G$ alone, including the family $\Qphi$ and the Casimirs of section~\ref{sec:invariants}.

We write $\Gamma_m=2b_m\langle g_mg_{m+1}\rangle$, so that $\sum_m\Gamma_m=\langle\fe{G}{LG}\rangle$; up to normalisation, $\Gamma_m$ is the contribution of the Hermite pair $(m,m+1)$ to $\Gamma$ in CMM's Hermite form (B32).

\begin{proposition}[No Fj{\o}rtoft-type constraint]\label{prop:nofjortoft}
Let $\Lam\notin[0,1]$ and consider one compressive branch of collisionless KRMHD in a periodic box.
\begin{itemize}
\item[(a)] If $\fe{e_m}{\varphi(L)e_n}=0$ for all $m\ne n$, for a real function $\varphi$ such that the $e_n$ lie in the domain of $\varphi(L)$, then $\varphi$ is constant almost everywhere on $\mathbb R$ (at the nodes $u_j$ in the truncated system). The only members of the family $\Qphi$ whose density in $m$ is diagonal are the multiples of $Q_1$.
\item[(b)] A positive invariant can have a Hermite weight that grows with $m$: $Q_{u^2}=\int\fe{LG}{LG}\,\dd^3r\ge0$, and $Q_{u^2}(e_m)/Q_1(e_m)=m+\frac12$ for every $m\ge2$.
\item[(c)] Let $\mathcal I$ be a functional of $G$ that is defined and constant along every solution of the branch equation with $\Phi=\Psi=0$ that starts from a state with a finite Hermite expansion in $u$, in particular any invariant of collisionless KRMHD that depends on $G$ alone and is defined along these solutions, and let $\mathcal I=\sum_mE_{\mathcal I}(m)$ be any decomposition of it over Hermite indices. There are no constants $c>0$ and $n>0$ such that $E_{\mathcal I}(m)\ge c\,m^nW(m)$ for every $m$ and every state, and none such that $E_{\mathcal I}(m)\ge c\,m^n|\Gamma_m|$ for every $m$ and every state.
\item[(d)] $\Gamma$ and $\Hph$ are odd under the mirror map \eqref{eq:mirror}. Every mirror-symmetric state has $\Gamma_m=0$ for all $m$ and $\Hph=0$, whatever its $W(m)$.
\end{itemize}
\end{proposition}

\begin{proof}
(a) The hypothesis is needed only for $n=0$. If $\fe{e_m}{\varphi(L)e_0}=0$ for every $m\ne0$, then, because the $e_m$ are complete in $(\mathcal W,\fe{\cdot}{\cdot})$ and $\fe{e_m}{e_0}=p_0\delta_{m0}$, $\varphi(L)e_0=\theta\,e_0$ with $\theta=\fe{e_0}{\varphi(L)e_0}/p_0$. The slice amplitude of $e_0$ is $p_0\sqrt{\varrho(u)}$, which is positive for every $u$, and that of $\varphi(L)e_0$ is $\varphi(u)\,p_0\sqrt{\varrho(u)}$ (section~\ref{sec:invariants}). Hence $\varphi=\theta$ almost everywhere. In the truncated system the slice components of $e_0$ are $p_0R_{0j}=p_0/\sqrt{n_j}\neq0$, and the same argument, applied to the matrix $\Pw\varphi(A)$, gives $\varphi(u_j)=\theta$ at every node.

(b) By the symmetry of $L$, $\fe{G}{L^2G}=\fe{LG}{LG}$, which is non-negative because $\fe{\cdot}{\cdot}$ is positive for $\Lam\notin[0,1]$. For $m\ge2$, $Le_m=\sqrt{m/2}\,e_{m-1}+\sqrt{(m+1)/2}\,e_{m+1}$ and $p_{m\pm1}=1$, so $\fe{Le_m}{Le_m}=m/2+(m+1)/2$, while $\fe{e_m}{e_m}=1$. (At $\Lam=\infty$ the ratio is $m+\frac12$ for every $m$.)

(c) The fields $\Phi=\Psi=0$ solve the RMHD equations \eqref{eq:rmhd}, and the slice equation \eqref{eq:slice} reduces to $\partial_ta_u+su\,\partial_za_u=0$. Each slice translates along $z$, so every $\Qphi$ and every Casimir $\int\mathcal F(a_u)\,\dd^3r$ \eqref{eq:casimir} is conserved. Take initial data with a finite Hermite expansion in $u$ and smooth spatial dependence, with $\langle\fe{G}{G}\rangle>0$ and no Fourier component at $k_z=0$. By hypothesis $\mathcal I$ is defined and constant along this solution, and $G$ lies in the domain of $L$. By \eqref{dec:gn}, each Fourier mode has $g_m(\boldsymbol k,t)=\int\sqrt\varrho\,r_m\,a_u(\boldsymbol k,0)\,e^{-\ii sk_zut}\,\dd u$. The function $\sqrt\varrho\,r_m\,a_u(\boldsymbol k,0)$ is integrable in $u$, by the Cauchy--Schwarz inequality, so $g_m(\boldsymbol k,t)\to0$ as $t\to\infty$ by the Riemann--Lebesgue lemma. The bound $|g_m(\boldsymbol k,t)|^2\le p_m^{-1}\int|a_u(\boldsymbol k,0)|^2\,\dd u$ is summable over $\boldsymbol k$, so by dominated convergence $W(m,t)\to0$ and $\Gamma_m(t)\to0$ for every $m$. Suppose now that $E_{\mathcal I}(m)\ge c\,m^nW(m)$ for every $m$ and every state. Then $E_{\mathcal I}(m)\ge0$, and for every $m_*$, $\mathcal I\ge c\,m_*^n\sum_{m>m_*}W(m,t)$ at all times. Since $\sum_mW(m,t)=\langle\fe{G}{G}\rangle$ is constant, the right-hand side tends to $c\,m_*^n\langle\fe{G}{G}\rangle$, which exceeds the constant $\mathcal I$ for $m_*$ large enough. If instead $E_{\mathcal I}(m)\ge c\,m^n|\Gamma_m|$, take data with $\langle\fe{G}{LG}\rangle\ne0$. The series $\sum_m\Gamma_m(t)$ then converges absolutely, and since $G(t)$ stays in the domain of $L$, its sum is the conserved $\langle\fe{G}{LG}\rangle$. Hence $\sum_{m>m_*}|\Gamma_m(t)|\ge|\langle\fe{G}{LG}\rangle-\sum_{m\le m_*}\Gamma_m(t)|$, which tends to $|\langle\fe{G}{LG}\rangle|$, and the same argument applies.

(d) By the parity rule of section~\ref{sec:symmetry}, $\langle g_mg_{m+1}\rangle$ is mirror-odd for every $m$, and so is $\Hph$. A mirror-symmetric state therefore has $\Gamma_m=-\Gamma_m=0$ and $\Hph=0$, while $W(m)$, which is mirror-even, can take any value.
\end{proof}

Linear phase mixing in a uniform field conserves every member of the family and every Casimir exactly, while it moves all of the free energy, and all of $\Gamma$, to arbitrarily large $m$. Whatever $m$-density is assigned to an invariant, that density therefore cannot satisfy a Fj{\o}rtoft-type inequality with $W(m)$ or $|\Gamma_m|$. Free streaming from density data at a single wavenumber shows how a positive invariant with a weight that grows with $m$ accommodates this. By \eqref{appE:poisson}, the fraction of the free energy in the moments $m\le m_*$ is the probability that a Poisson variable of mean $\lambda=(sk_zt)^2/2$ does not exceed $m_*$. Per unit free energy, the part of $Q_{u^2}$ diagonal in $m$ grows as $\lambda+\frac12$, and the off-diagonal part, $-\lambda$, keeps $Q_{u^2}$ constant. In two-dimensional hydrodynamics a solution of this kind is impossible, because the energy above a wavenumber $k_*$ is bounded by the enstrophy divided by $k_*^2$.

In a truncated hierarchy the same holds until the recurrence time. At $M=512$ the fraction of the free energy in $m\le10$ falls from 1 to $3\times10^{-29}$ by $sk_zt=14$ for free streaming, and to $3\times10^{-4}$ for $\Lam=\sqrt5$. For data with $\Gamma\ne0$, $g(0)=e_0+0.8\,e_1$ at $M=3000$, the fraction $\sum_{m\le20}|\Gamma_m|/|\sum_m\Gamma_m|$ falls to $4\times10^{-28}$ for free streaming and to $4\times10^{-18}$ for $\Lam=\sqrt5$ by $sk_zt=30$, while $\sum_mm\,W(m)$ grows from 0.6 to 739 and to 530.

The reason lies in the slice structure. Every invariant of section~\ref{sec:invariants} is a functional of single-slice statistics: $\Qphi$ depends only on the diagonal $C(u,u)$, a Casimir $\int\mathcal F(a_u)\,\dd^3r$ only on the one-point distribution of a single slice, and a cross-sector invariant \eqref{eq:crosssector} only on slices of different sectors at the same $u$. The Hermite transfer \eqref{eq:fluxCD} is $-\iint K_m\,\partial_tC$, a change in the correlation between different slices of one sector, with $\partial_tC(u,u)=0$. The invariants do not depend on these correlations, so their conservation places no constraint on the transfer. The expectation of \citet{Chandran2026} that $\Gamma$, like $W$, is transferred to large $m$ is therefore borne out by the exact solutions considered here, for a reason that does not rely on the sign-indefiniteness of $\Gamma$. The argument concerns exact transient solutions of collisionless KRMHD, in the limit $k_\perp\rho_i\to0$. The criterion of \citet{Chandran2026} refers to inertial-range fluxes in a forced and dissipative steady state, and a steady-state version of the argument is not attempted here. At finite $k_\perp\rho_i$ the velocity-space operator depends on the wavenumber (section~\ref{sec:invariants}), and the argument would have to be redone.
